\section{A Generalization of Alternating Multilinear Forms}

Let \(\mathrm{V}\) be a right vector space over the field \(\mathrm{D}\) , of finite dimension \(n \geqslant 0\) . We denote by \(\mathrm{B(V)}\) the set of bases of \(\mathrm{V}\) and by \(\Omega(\mathrm{V})\) the set of mappings \(\omega\) from \(\mathrm{B(V)}\) to \(\mathrm{D}_{\mathrm{ab}}^{*}\) that satisfy the following two conditions:

\begin{enumerate}
\def\labelenumi{\alph{enumi})}
\tightlist
\item
  If \(\lambda_1, \ldots, \lambda_n\) are the elements of \(D^*\) , then we have
\end{enumerate}

\[
\omega (v _ {1} \lambda_ {1}, \ldots , v _ {n} \lambda_ {n}) = \pi (\lambda_ {1} \dots \lambda_ {n}) \omega (v _ {1}, \ldots , v _ {n})\tag{1}
\]

for every basis \((v_{1},\ldots ,v_{n})\) of V.

\begin{enumerate}
\def\labelenumi{\alph{enumi})}
\setcounter{enumi}{1}
\tightlist
\item
  If \(i\) and \(j\) are two distinct integers in the interval \([1, n]\) , then we have
\end{enumerate}

\[
\omega (v _ {1}, \ldots , v _ {i - 1}, v _ {i} + v _ {j}, v _ {i + 1}, \ldots , v _ {n}) = \omega (v _ {1}, \ldots , v _ {n})\tag{2}
\]

for every basis \((v_{1},\ldots ,v_{n})\) of V.

Let \(\omega\) be an element of \(\Omega(\mathrm{V})\) . Let \((v_1, \ldots, v_n)\) be a basis of \(\mathrm{V}\) and \(i\) an integer in the interval \([1, n-1]\) ; by property b), we have

\[
\begin{array}{r l} & {\omega (v _ {1}, \ldots , v _ {i}, v _ {i + 1}, \ldots , v _ {n}) = \omega (v _ {1}, \ldots , v _ {i} + v _ {i + 1}, v _ {i + 1}, \ldots , v _ {n})} \\ & {\qquad = \omega (v _ {1}, \ldots , v _ {i} + v _ {i + 1}, v _ {i + 1} - (v _ {i} + v _ {i + 1}), \ldots , v _ {n})} \\ & {\qquad = \omega (v _ {1}, \ldots , v _ {i} + v _ {i + 1}, - v _ {i}, \ldots , v _ {n})} \\ & {\qquad = \omega (v _ {1}, \ldots , (v _ {i} + v _ {i + 1}) - v _ {i}, - v _ {i}, \ldots , v _ {n}),} \end{array}
\]

and therefore

\[
\omega (v _ {1}, \dots , v _ {i}, v _ {i + 1}, \dots , v _ {n}) = \pi (- 1) \omega (v _ {1}, \dots , v _ {i + 1}, v _ {i}, \dots , v _ {n}).\tag{3}
\]

Since the symmetric group \(\mathfrak{S}_n\) is generated by the transpositions of two consecutive elements of the interval \([1,n]\) (I, §5, No.~7, p.~63, Proposition 9), formula (3) generalizes to

\[
\omega (v _ {\sigma (1)}, \ldots , v _ {\sigma (n)}) = \pi (\varepsilon (\sigma)) \omega (v _ {1}, \ldots , v _ {n}),\tag{4}
\]

where \(\sigma\) belongs to \(\mathfrak{S}_n\) and \(\varepsilon (\sigma)\) is its signature (I, §5, No.~7, p.~64).

Formula (2) generalizes as follows. First, for every \(\lambda\) in \(\mathrm{D}^*\) , we have

\[
\begin{array}{r l} & {\omega (v _ {1}, \ldots , v _ {n}) = \pi (\lambda) \omega (v _ {1}, \ldots , v _ {i} \lambda^ {- 1}, \ldots , v _ {n})} \\ & {\qquad = \pi (\lambda) \omega (v _ {1}, \ldots , v _ {i} \lambda^ {- 1} + v _ {j}, \ldots , v _ {n})} \\ & {\qquad = \omega (v _ {1}, \ldots , (v _ {i} \lambda^ {- 1} + v _ {j}) \lambda , \ldots , v _ {n}),} \end{array}
\]

that is,

\[
\omega (v _ {1}, \dots , v _ {n}) = \omega (v _ {1}, \dots , v _ {i} + v _ {j} \lambda , \dots , v _ {n}).\tag{5}
\]

By an immediate induction, we deduce

\[
\omega (v _ {1}, \ldots , v _ {n}) = \omega (v _ {1}, \ldots , v _ {i} + w, \ldots , v _ {n})\tag{6}
\]

for every linear combination \(w = \sum_{j \neq i} v_j \lambda_j\) of the vectors \(v_j\) for j different from i in the interval \([1, n]\) .

